\chapter{电路与安全}\label{ch:g6:circuits-and-safety}

三年来，你已经学会了连接环路、灯泡和支路，成为一名
电路的搭建者。今年，你将成为一种更稀有的角色：
电路的\emph{书写者}。各国的电工读写和绘制的是同一套
无需文字的语言——这套语法，连同一章来之不易的
安全经验，正是区分巧妙接线与
烧伤手指的关键。

\section{无需文字的语言}

\begin{definition}[电路图]\label{def:g6:circuits-and-safety:diagram}
\emph{电路图}\index{电路图}是电路的标准画法：
每个元件都用约定的
\emph{符号}\index{符号}表示，导线用带直角的直线画出。
电路图表明电路是如何
\emph{连接}的——谁给谁供电、何处分叉——
而刻意忽略实际导线在桌面上如何蜿蜒。
\end{definition}

\begin{omfigure}
\begin{tikzpicture}[scale=0.8]
  \draw (0.4,3.4) to[battery1] (2.4,3.4);
  \node[right, font=\small] at (2.7,3.4)
    {电池（长线为 $+$ 极）};
  \draw (0.4,2.2) to[lamp] (2.4,2.2);
  \node[right, font=\small] at (2.7,2.2) {灯泡};
  \draw (0.4,1.0) to[spst] (2.4,1.0);
  \node[right, font=\small] at (2.7,1.0) {开关（断开）};
  \draw (7.4,3.4) to[cspst] (9.4,3.4);
  \node[right, font=\small] at (9.7,3.4) {开关（闭合）};
  \draw (7.4,2.2) -- (8.05,2.2);
  \draw[thick] (8.4,2.2) circle (0.35);
  \node[font=\small] at (8.4,2.2) {M};
  \draw (8.75,2.2) -- (9.4,2.2);
  \node[right, font=\small] at (9.7,2.2) {电动机};
  \draw (7.4,1.0) -- (8.05,1.0);
  \draw[thick] (8.4,1.0) circle (0.35);
  \draw[thick] (8.16,0.76) -- (8.64,1.24);
  \draw (8.75,1.0) -- (9.4,1.0);
  \node[right, font=\small] at (9.7,1.0) {蜂鸣器};
\end{tikzpicture}

{\small 符号字母表：六种符号几乎可以画出本书的
每一个电路。这些符号在世界各地的
工坊中被同样地识读。}
\end{omfigure}

\begin{method}[从桌面到纸面，再从纸面回到桌面]\label{met:g6:circuits-and-safety:translate}
要\emph{画出}已搭好的电路：
\begin{enumerate}
  \item 列出各个元件，并选定每个元件的符号；
  \item 用手指沿实际回路，从电池的 $+$ 极出发，
        记下各元件的先后顺序和每一处分叉；
  \item 重画成整齐的矩形：符号间距适中，导线
        平直，转角为直角——连接关系不变，导线不得弯曲。
\end{enumerate}
要根据电路图\emph{搭建}电路，则把上述过程反向进行——
搭建完成后，照常检查：从 $+$ 极出发追踪的每一个
回路，都必须经过至少一个元件回到 $-$ 极。
\end{method}

\begin{example}[同一电路的两幅画像]\label{ex:g6:circuits-and-safety:portraits}
桌面上有：一节电池，一团绕过铅笔盒的鳄鱼夹导线，
一盏靠在书上的灯泡，以及一个悬挂着的开关。
纸上则是一个整洁的矩形：电池在左，开关在上，
灯泡在右。乱线与矩形是\emph{同一个
电路}：每一处连接都完全相同。要判断两个
电路是否相同，只需问一个问题：
谁和谁相连？
\end{example}

\begin{example}[电动机与蜂鸣器登场]\label{ex:g6:circuits-and-safety:motor}
套件中的两位新成员，在电路图中的表现与灯泡完全相同：
\emph{电动机}\index{电动机}（符号为圆圈中的 M）只要有电流通过就转动——
例如风扇、玩具车、电钻；
\emph{蜂鸣器}\index{蜂鸣器}（画成带敲击标记的圆）只要有电输入便发声——
例如门铃、警报器。串联和并联的规则对它们同样适用，
不做改变：串联的开关可以控制它们；并联的各条支路
使它们彼此独立。
\end{example}

\section{新语言中的老规则}

\begin{example}[像行家一样读图]\label{ex:g6:circuits-and-safety:reading}
对任意电路图，可用并联电路一章中的“手指行走”法检验：
从 $+$ 极出发，走遍每一条通往 $-$ 极的路线。全部
位于同一条路线上的元件：串联——任一处断开，全部停止，
顺序任意，亮度共享。在圆点节点处分叉的路线：并联——
各条支路彼此独立，每条亮度都满，开关位置因此成为
需要真正权衡的问题。语言变了，规律没变。
\end{example}

\begin{omfigure}
\begin{tikzpicture}[scale=0.8]
  \draw (0,0) to[battery1] (0,3.2)
    to[cspst] (2.6,3.2);
  \draw (2.6,3.2) to[short] (2.6,2.4)
    to[lamp] (5.8,2.4) to[short] (5.8,3.2);
  \draw (2.6,3.2) to[short] (2.6,4.0) -- (3.4,4.0);
  \draw[thick] (3.75,4.0) circle (0.35);
  \node[font=\small] at (3.75,4.0) {M};
  \draw (4.1,4.0) -- (5.8,4.0) to[short] (5.8,3.2);
  \draw (5.8,3.2) to[short] (7.4,3.2) to[short] (7.4,0)
    to[short] (0,0);
  \fill (2.6,3.2) circle (2pt);
  \fill (5.8,3.2) circle (2pt);
\end{tikzpicture}

{\small 搭建之前先读图：公共路段上有一个开关，
然后分叉——一条支路接灯泡，另一条支路接风扇电动机。
开关同时控制两者；两条支路互不受对方故障影响。}
\end{omfigure}

\section{伤痕之章}

用电安全规则是用火灾乃至更惨重的代价换来的。这里
列出了这些规则及其理由——行家两者都懂。

\begin{definition}[短路]\label{def:g6:circuits-and-safety:short}
\emph{短路}\index{短路}发生于一段裸露的
导电通路将电池（或插座）的两极
\emph{直接}连通，跳过每一个用电器。
途中没有任何有用的东西阻挡，电流便疯狂冲过——导线发热，绝缘层熔化，
电池报废，市电线路还可能起火。下一年的电路章将深入研究这个元凶；
今年则先认清它的面目，
并把它挡在你的搭建之外。
\end{definition}

\begin{remark}[家中的守卫]\label{rem:g6:circuits-and-safety:guards}
打开家里的配电箱看看：一排排小开关——
\emph{断路器}——家中大型并联电路的每条支路各一个。
每个断路器守护着自己的支路：一旦出现故障，让电流猛冲——
比如短路、插座过载——
守卫便瞬间自行跳开，在导线过热之前切断该支路。
断路器“跳闸”时，它并非
出了毛病，而是刚刚完成了自己唯一的职责。
复位之前，务必先找出故障。
\end{remark}

\begin{remark}[规则及其理由]\label{rem:g6:circuits-and-safety:rules}
\begin{enumerate}
  \item \emph{只用电池做实验，插座绝不能碰。}
        市电的电压要高出数百倍——
        隔着潮湿的皮肤，它就可能让心脏停跳。
  \item \emph{不要用裸导线直接横跨电池两极。}这就是
        一次短路：电池报废、导线发热、手指烫伤。
  \item \emph{水与市电绝不能共处于
        房内伸手可及之处：}浴缸旁不放电器，湿手不碰插头——
        潮湿的皮肤就是导体。
  \item \emph{绝不要把任何东西捅进插座。}配电箱里的守卫
        动作虽快，但你绝不能拿它来赌。
  \item \emph{破损的电线要扔进垃圾桶，}而不是塞进抽屉：
        开裂的绝缘层就像一道留有缺口的栅栏。
\end{enumerate}
\end{remark}

\section{练习}

\begin{exercise}[$\star$]\label{exo:g6:circuits-and-safety:1}
电路图忠实地展示了什么，
又刻意忽略了什么？
\end{exercise}

\begin{exercise}[$\star$]\label{exo:g6:circuits-and-safety:2}
默写下列符号：电池；\ 灯泡；\ 断开的开关；\
闭合的开关；\ 电动机；\ 蜂鸣器。哪个标记表示电池的
$+$ 极？
\end{exercise}

\begin{exercise}[$\star$]\label{exo:g6:circuits-and-safety:3}
画出门铃的电路图：电池、按钮开关和蜂鸣器
接成一个回路。手指一松开，门铃为何停止响？
\end{exercise}

\begin{exercise}[$\star$]\label{exo:g6:circuits-and-safety:4}
用同一盒元件搭出了两个电路，它们算不算
“同一个电路”？由哪一个问题即可判定？
\end{exercise}

\begin{exercise}[$\star$]\label{exo:g6:circuits-and-safety:5}
在本章的灯泡与电动机电路图中，若灯泡烧坏，风扇会怎样？
若开关断开，两者又会如何？哪些旧规则
可以用来解答？
\end{exercise}

\begin{exercise}[$\star$]\label{exo:g6:circuits-and-safety:6}
什么是短路？它为什么会使导线发热？请列举其中
两个受害对象。
\end{exercise}

\begin{exercise}[$\star$]\label{exo:g6:circuits-and-safety:7}
配电箱里的一个断路器跳闸了。它刚刚为住宅做了
什么？复位之前应当先做什么？
\end{exercise}

\begin{exercise}[$\star$]\label{exo:g6:circuits-and-safety:8}
说明每条规则背后的理由：浴缸附近不放电器；\
破损电线要扔进垃圾桶；\ 玩电池游戏可以，插座绝对不行。
\end{exercise}

\begin{exercise}[$\star\star$]\label{exo:g6:circuits-and-safety:9}
将以下描述转换为电路图（用文字或图画均可）：一个电池；由其
$+$ 极引出一根导线接到闭合的开关；由开关分出一个分叉——一条
支路经过蜂鸣器，一条经过灯泡；各支路重新汇合并回到 $-$ 极。
预测能听到什么、看到什么，以及若
只给蜂鸣器所在的支路再单独加装一个开关并断开，会有什么变化？
\end{exercise}

\begin{exercise}[$\star\star$]\label{exo:g6:circuits-and-safety:10}
一台玩具风扇的接线方式如下：电池 $+$ 极接电动机，电动机接电池 $-$ 极，
此外——“为了牢固”——再加一根裸导线，直接从
$+$ 极连到 $-$ 极。风扇几乎转不动，而那根多出来的导线越来越烫。
请运用\cref{def:g6:circuits-and-safety:short}诊断：电流更倾向于走哪条
路径，这种接法为什么危险？
\end{exercise}

\begin{exercise}[$\star\star$]\label{exo:g6:circuits-and-safety:11}
爷爷在车间里有个窍门：碰任何灯具之前，他
先关掉墙上的开关，\emph{并且}关掉
配电箱里该支路的断路器。运用串联电路的思路，说明每一次
切断各自达到了什么效果——以及这位谨慎的人为什么两步都做。
\end{exercise}

\begin{exercise}[$\star\star\star$]\label{exo:g6:circuits-and-safety:12}
设计一个木偶剧场的布线：一个总开关，能
让全部灯暗下来；两盏舞台灯必须彼此独立失效；
一台幕布电动机和一个开演蜂鸣器，各自带一个控制开关。
请描述（或画出）电路图——哪些用电器走公共干路，
哪些走支路，每个开关装在何处——并
对照你掌握的串联与并联的每条规则
检验你的设计。
\end{exercise}

\section{问题：给花园小屋布线}

\begin{problem}[{周末问题——一家人给自家花园小屋布线；
先画图，后动螺丝刀；检查员的
检查清单}]\label{pb:g6:circuits-and-safety:1}
这家人先在纸上规划小屋里这套电池供电的小电路——
就像专业人员那样。

\medskip\noindent\textbf{Part I——纸上的方案。}
需求清单：一盏吊灯在门口用开关控制；一盏工作台灯
带独立开关；一台小通风扇（电动机）带独立
开关。
\begin{enumerate}
  \item 为什么三个用电器必须走并联支路，
        而不能共用一条串联回路？请举出两种
        若串联就会失效的便利。
  \item 每条支路各有一个开关。每个开关控制什么，
        而所有开关都无法控制什么？
  \item 总开关应装在哪里，才能让整个小屋
        同时关灯、断电？
  \item 画出（或精确描述）完整电路图：电池、
        总开关、三条支路，共四个开关。
\end{enumerate}

\medskip\noindent\textbf{Part II——安装后的检查。}
\begin{enumerate}[resume]
  \item 装好并合上开关后，吊灯确实亮了——但
        很暗；工作台灯只有在
        \emph{风扇}的开关闭合时才亮。顺着线路逐段排查，发现
        工作台灯与风扇共用一条支路，二者
        一前一后相连。说明它们的连接方式，并用旧规则解释
        这两种现象。
  \item 只用一根导线修复，使电路恢复原方案。
  \item 整理时，一把螺丝刀短暂地搭在
        电池两个裸露的接线柱上，螺丝刀随即发热。指出
        这是什么现象（短路），并说明事后家人应检查
        电池的哪些情况。
  \item 检查员家长又加一条规则：“每个接头都包好，
        任何地方都不许有裸露的铜线。”这一条规则防范的是
        哪两种危险？
\end{enumerate}

\medskip\noindent\textbf{Part III——暴风雨之夜。}
\begin{enumerate}[resume]
  \item 一个雨夜：小屋屋顶漏水，水滴到工作台上。
        为什么即使在电池供电的小屋里，家人也不能
        用湿手碰电线——同样的接触在市电供电的房子里，
        后果为何要重得多？
  \item 房子里的灯忽明忽暗，一个断路器跳闸。
        断路器检测到了什么，
        以及明智的处理顺序是什么？
  \item 早上，全家争论要不要加装一个门铃——蜂鸣器
        在小屋，按钮开关在房子里，两者之间用一根长的双芯导线
        相连。它与小屋的电路是串联还是并联？
        它的电流从哪里来？画出这个回路的草图。
  \item 为小屋写一份全家的三行安全守则，
        每行一条，分别针对插座、水和裸露的电线。
\end{enumerate}
\end{problem}
